\section{Installation on the Raspberry Pi}

\subsection{System Requirements}

Automated image acquisition requires a Raspberry Pi 4 or newer, a Raspberry Pi
Camera Module v3, Raspberry Pi OS Bookworm or newer, and Picamera2. The
scheduler is installed as a systemd service, so the Raspberry Pi should be set
up with a working network connection, correct system time, and enough storage
space for the planned experiment.

The installation below assumes that Phenopi will run under a dedicated user
account named \texttt{phenopi}. This keeps the acquisition service separate
from personal user accounts and makes file permissions easier to manage.

\subsection{Expected Installation Layout}

The provided service file assumes the following installation layout:

\begin{verbatim}
  /home/phenopi/phenotyping
  /home/phenopi/venvs/phenopi
\end{verbatim}

The Phenopi repository is stored in \texttt{/home/phenopi/phenotyping}. The
Python virtual environment is stored in \texttt{/home/phenopi/venvs/phenopi}.
If different paths are used, ensure to update the paths in
\texttt{deploy/phenopi-scheduler.service} also.

\subsection{Creating the Phenopi User}

Create a dedicated user for running the scheduler.

\begin{verbatim}
sudo useradd -m -s /bin/bash phenopi
\end{verbatim}

Create the directory that will contain the Python virtual environment.

\begin{verbatim}
sudo -u phenopi mkdir -p /home/phenopi/venvs
\end{verbatim}

\subsection{Installing Phenopi}

Clone the Phenopi repository into the expected location.

\begin{verbatim}
sudo -u phenopi git clone https://github.com/kfreinders/phenopi.git 
/home/phenopi/phenotyping
\end{verbatim}

Install Picamera2 using the Raspberry Pi package manager.

\begin{verbatim}
sudo apt update
sudo apt install -y python3-picamera2
\end{verbatim}

Create the Python virtual environment. The \texttt{--system-site-packages}
option makes packages installed through \texttt{apt}, such as Picamera2,
available inside the virtual environment.

\begin{verbatim}
sudo -u phenopi python3 -m venv --system-site-packages 
/home/phenopi/venvs/phenopi
\end{verbatim}

Install the Python dependencies listed by the project.

\begin{verbatim}
sudo -u phenopi /home/phenopi/venvs/phenopi/bin/pip install -r 
/home/phenopi/phenotyping/requirements.txt
\end{verbatim}

\subsection{Checking the Scheduler Command}

Before installing the service, check that the scheduler can be started from the
virtual environment.

\begin{verbatim}
cd /home/phenopi/phenotyping
sudo -u phenopi /home/phenopi/venvs/phenopi/bin/python \
-m scripts.scheduling.run_scheduler --help
\end{verbatim}

A successful installation should display the command-line help page. If the
command fails, check the repository path, virtual environment path, installed
dependencies, and scheduler module name.

\subsection{Installing the Scheduler Service}

Copy the provided service file to the systemd service directory and reload
systemd so that it detects the new service.

\begin{verbatim}
sudo cp /home/phenopi/phenotyping/deploy/phenopi-scheduler.service \
/etc/systemd/system/phenopi-scheduler.service
sudo systemctl daemon-reload
\end{verbatim}

Enable and start the service so that it starts automatically after boot.

\begin{verbatim}
sudo systemctl enable phenopi-scheduler.service
sudo systemctl start phenopi-scheduler.service
\end{verbatim}

\subsection{Checking the Service}

Check whether the scheduler service is running.

\begin{verbatim}
sudo systemctl status phenopi-scheduler.service
\end{verbatim}

The service logs can be viewed with \texttt{journalctl}.

\begin{verbatim}
sudo journalctl -u phenopi-scheduler.service -f
\end{verbatim}

If the service fails to start, first check that the paths in
\texttt{deploy/phenopi-scheduler.service} match the installation paths. In
particular, verify \texttt{WorkingDirectory}, \texttt{ExecStart}, and whether
the \texttt{phenopi} user has permission to read the repository and write to
the configured output directories.
